Maxwell--Dirac 场方程已经固定局域源项与真空中的因果传播。介质、边界和开放几何并不改变 Maxwell 方程的局域来源，但会改变“给定电流以后场如何传播”的逆问题。在线性时不变环境中，这个逆问题由频域 Maxwell 算符及其推迟 Green 张量统一描述；散射、耗散和后续宏观 QED 都从同一个逆核出发。

先考虑空间局域但允许各向异性和色散的线性介质。相对介电率与相对磁导率张量分别记为 $\boldsymbol\varepsilon_r(\mathbf r,\omega)$、$\boldsymbol\mu_r(\mathbf r,\omega)$，二者均无量纲；$\omega$ 是角频率，单位为 $\mathrm{s^{-1}}$。采用全文时间 Fourier 约定 $\ee^{-\ii\omega t}$。消去磁场以后，电场满足的二阶微分算符定义为
\begin{equation}
 \boxed{
 \mathsf L_{\rm EM}(\omega)
 \equiv
 \nabla\times\boldsymbol\mu_r^{-1}(\mathbf r,\omega)\nabla\times
 -\frac{\omega^2}{c^2}\boldsymbol\varepsilon_r(\mathbf r,\omega).}
 \label{eq:Maxwell-operator-SI-r12}
\end{equation}
$\mathsf L_{\rm EM}$ 的量纲为长度的负二次方。材料本构、几何和边界条件都进入这个算符；若材料响应具有空间非局域性，$\boldsymbol\varepsilon_r$ 与 $\boldsymbol\mu_r$ 应相应提升为积分核，而同一算符结构仍然成立。

给定辐射出射条件，推迟并矢 Green 张量 $\mathbf G^R(\mathbf r,\mathbf r';\omega)$ 定义为 $\mathsf L_{\rm EM}$ 的逆核：
\begin{equation}
 \boxed{
 \mathsf L_{\rm EM}(\omega)
 \mathbf G^R(\mathbf r,\mathbf r';\omega)
 =\mathbf I\,\delta^{(3)}(\mathbf r-\mathbf r').}
 \label{eq:GR-Maxwell-SI-r12}
\end{equation}
其中 $\mathbf I$ 是三维单位并矢，$[\mathbf G^R]=\mathrm{m^{-1}}$。上标 $R$ 表示推迟边界条件：源在给定频率产生的场必须满足向外辐射而不是从无穷远流入的条件。因而 Green 张量不仅包含材料参数，还包含开放系统的物理边界条件。

若时域外加电流密度为 $\mathbf J(\mathbf r,t)$，则按全文 Fourier convention 定义的频谱电流 $\widetilde{\mathbf J}(\mathbf r,\omega)$ 的单位为 $\mathrm{A\,s\,m^{-2}}$。频域 Maxwell 方程可压缩成
\begin{equation}
\boxed{
\mathsf L_{\rm EM}(\omega)\widetilde{\mathbf E}(\mathbf r,\omega)
=\ii\mu_0\omega\widetilde{\mathbf J}(\mathbf r,\omega).}
\label{eq:maxwell-source-operator-dp68}
\end{equation}
两条旋度方程合并后得到与具体几何无关的线性源方程。利用\eqnrefs{eq:GR-Maxwell-SI-r12} 的推迟逆核，得到电流到电场的线性响应母式
\begin{equation}
 \boxed{
 \widetilde{\mathbf E}(\mathbf r,\omega)
 =\ii\mu_0\omega
 \int\dd^3 r'\,
 \mathbf G^R(\mathbf r,\mathbf r';\omega)
 \cdot\widetilde{\mathbf J}(\mathbf r',\omega).}
 \label{eq:E-from-GJ-SI-r12}
\end{equation}
量纲检查给 $[\mu_0\omega G^R\widetilde J\,\dd^3 r']=\mathrm{V\,s\,m^{-1}}$，与时间 Fourier 电场 $\widetilde{\mathbf E}$ 的单位一致。这个公式是后续 EELS、辐射和环境诱导自能的经典响应起点；具体电子轨迹只会改变源 $\widetilde{\mathbf J}$，不会改变环境核 $\mathbf G^R$。

入射场与结构散射可由参考算符分解。选取参考算符 $\mathsf L_0$ 及其推迟 Green 张量 $\mathbf G_0^R=\mathsf L_0^{-1}$，并定义散射势算符
\begin{equation*}
 \mathsf V_{\rm EM}\equiv\mathsf L_0-\mathsf L_{\rm EM}.
\end{equation*}
同一源在真实环境中的场满足
\begin{equation}
\boxed{
\widetilde{\mathbf E}
=\widetilde{\mathbf E}_{\rm inc}+\mathbf G_0^R\mathsf V_{\rm EM}\widetilde{\mathbf E}.}
\label{eq:maxwell-lippmann-schwinger-dp68}
\end{equation}
右边第二项表示场先由结构散射一次，再由参考环境传播；把其中的总场再次用同一方程代回，就自动产生二次、三次直到任意次数的散射。只要相应预解式存在，这个级数可重求和为电磁 $T$ 算符
\begin{equation}
 \boxed{
 \mathsf T_{\rm EM}
 =\mathsf V_{\rm EM}
 (\mathbf I-\mathbf G_0^R\mathsf V_{\rm EM})^{-1}.}
 \label{eq:maxwell-t-operator-def-r16}
\end{equation}
于是频域总场 $\widetilde{\mathbf E}=\widetilde{\mathbf E}_{\rm inc}+\widetilde{\mathbf E}_{\rm sc}$ 的散射部分为
\begin{equation}
 \boxed{
 \widetilde{\mathbf E}_{\rm sc}
 =\mathbf G_0^R\mathsf T_{\rm EM}\widetilde{\mathbf E}_{\rm inc}.}
 \label{eq:maxwell_scattered_field_t_operator}
\end{equation}
这里 $\widetilde{\mathbf E}_{\rm inc}$ 是同一外部源在参考环境 $\mathsf L_0$ 中产生的 Fourier 电场。球、圆柱、平面界面和数值离散结构的差别，只在于如何表示并求解同一个 $\mathsf T_{\rm EM}$。

若局域介质满足互易本构条件
$\boldsymbol\varepsilon_r^T=\boldsymbol\varepsilon_r$ 与
$\boldsymbol\mu_r^T=\boldsymbol\mu_r$，并对两组场采用同一辐射边界条件，则分部积分后 Maxwell 算符满足
\begin{equation}
\boxed{
\int\dd^3 r\,\mathbf u\!\cdot\!\mathsf L_{\rm EM}\mathbf v
=\int\dd^3 r\,\mathbf v\!\cdot\!\mathsf L_{\rm EM}\mathbf u.}
\label{eq:maxwell-bilinear-reciprocity-dp68}
\end{equation}
这条式子只是说：在互易介质中，交换“源场”和“测试场”不会改变这一双线性响应。把两组测试场选成点源产生的 Green 场，就得到 Lorentz 互易关系
\begin{equation}
 \boxed{
 \mathbf G^R(\mathbf r,\mathbf r';\omega)
 =(\mathbf G^R)^T(\mathbf r',\mathbf r;\omega).}
 \label{eq:maxwell-green-reciprocity-r16}
\end{equation}
互易性约束交换源点与观测点时的张量转置；它与“无损”是不同条件，因此有吸收的互易介质仍可满足此式。

Green 张量还把材料耗散与场的谱密度联系起来。一般算符恒等式为
\begin{equation}
 \boxed{
 \mathbf G^R-(\mathbf G^R)^\dagger
 =(\mathbf G^R)^\dagger
 (\mathsf L_{\rm EM}^\dagger-\mathsf L_{\rm EM})
 \mathbf G^R.}
 \label{eq:resolvent-imag-identity-r12}
\end{equation}
左边是因果预解式的反厄米部分，右边则显示它只能来自 Maxwell 算符自身的反厄米部分，也就是吸收通道。若 $\boldsymbol\mu_r$ 为实厄米张量且耗散只来自电响应，则上式化为
\begin{equation}
 \boxed{
 \frac{\omega^2}{c^2}
 \int\dd^3 s\,
 \mathbf G^R(\mathbf r,\mathbf s;\omega)
 \cdot\operatorname{Im}\boldsymbol\varepsilon_r(\mathbf s,\omega)
 \cdot(\mathbf G^R)^\dagger(\mathbf r',\mathbf s;\omega)
 =\operatorname{Im}\mathbf G^R(\mathbf r,\mathbf r';\omega).}
 \label{eq:green-optical-identity-SI-r12}
\end{equation}
这条光学恒等式说明同一个 $\operatorname{Im}\mathbf G^R$ 同时记录传播与不可逆耗散。存在磁损耗时，右侧仍由\eqnrefs{eq:resolvent-imag-identity-r12} 控制，只需在耗散核中加入 $\operatorname{Im}\boldsymbol\mu_r^{-1}$ 的贡献。
